
\documentclass{article}
%%%%%%%%%%%%%%%%%%%%%%%%%%%%%%%%%%%%%%%%%%%%%%%%%%%%%%%%%%%%%%%%%%%%%%%%%%%%%%%%%%%%%%%%%%%%%%%%%%%%%%%%%%%%%%%%%%%%%%%%%%%%%%%%%%%%%%%%%%%%%%%%%%%%%%%%%%%%%%%%%%%%%%%%%%%%%%%%%%%%%%%%%%%%%%%%%%%%%%%%%%%%%%%%%%%%%%%%%%%%%%%%%%%%%%%%%%%%%%%%%%%%%%%%%%%%
\usepackage{amssymb}

%TCIDATA{OutputFilter=LATEX.DLL}
%TCIDATA{Version=5.50.0.2960}
%TCIDATA{<META NAME="SaveForMode" CONTENT="1">}
%TCIDATA{BibliographyScheme=Manual}
%TCIDATA{Created=Saturday, June 14, 2014 09:45:02}
%TCIDATA{LastRevised=Friday, June 27, 2014 09:42:46}
%TCIDATA{<META NAME="GraphicsSave" CONTENT="32">}
%TCIDATA{<META NAME="DocumentShell" CONTENT="Standard LaTeX\Blank - Standard LaTeX Article">}
%TCIDATA{CSTFile=40 LaTeX article.cst}

\newtheorem{theorem}{Theorem}
\newtheorem{acknowledgement}[theorem]{Acknowledgement}
\newtheorem{algorithm}[theorem]{Algorithm}
\newtheorem{axiom}[theorem]{Axiom}
\newtheorem{case}[theorem]{Case}
\newtheorem{claim}[theorem]{Claim}
\newtheorem{conclusion}[theorem]{Conclusion}
\newtheorem{condition}[theorem]{Condition}
\newtheorem{conjecture}[theorem]{Conjecture}
\newtheorem{corollary}[theorem]{Corollary}
\newtheorem{criterion}[theorem]{Criterion}
\newtheorem{definition}[theorem]{Definition}
\newtheorem{example}[theorem]{Example}
\newtheorem{exercise}[theorem]{Exercise}
\newtheorem{lemma}[theorem]{Lemma}
\newtheorem{notation}[theorem]{Notation}
\newtheorem{problem}[theorem]{Problem}
\newtheorem{proposition}[theorem]{Proposition}
\newtheorem{remark}[theorem]{Remark}
\newtheorem{solution}[theorem]{Solution}
\newtheorem{summary}[theorem]{Summary}
\newenvironment{proof}[1][Proof]{\noindent\textbf{#1.} }{\ \rule{0.5em}{0.5em}}
\input{tcilatex}
\begin{document}


\section{A Concrete Example}

On the space $\mathbb{C}^{2}$ the partial isometries $W_{1}=\left( 
\begin{array}{cc}
0 & 1 \\ 
0 & 0%
\end{array}%
\right) ,$ and $W_{2}=\left( 
\begin{array}{cc}
0 & 0 \\ 
1 & 0%
\end{array}%
\right) $ have the properties 
\begin{eqnarray}
W_{1}W_{1}^{\dagger } &=&\left( 
\begin{array}{cc}
0 & 1 \\ 
0 & 0%
\end{array}%
\right) \left( 
\begin{array}{cc}
0 & 0 \\ 
1 & 0%
\end{array}%
\right) =\left( 
\begin{array}{cc}
0 & 1 \\ 
0 & 0%
\end{array}%
\right) \left( 
\begin{array}{cc}
0 & 1 \\ 
0 & 0%
\end{array}%
\right) ^{\dagger }=\left( 
\begin{array}{cc}
1 & 0 \\ 
0 & 0%
\end{array}%
\right) =P_{1} \\
W_{2}W_{2}^{\dagger } &=&\left( 
\begin{array}{cc}
0 & 0 \\ 
1 & 0%
\end{array}%
\right) \left( 
\begin{array}{cc}
0 & 1 \\ 
0 & 0%
\end{array}%
\right) =\left( 
\begin{array}{cc}
0 & 1 \\ 
0 & 0%
\end{array}%
\right) ^{\dagger }\left( 
\begin{array}{cc}
0 & 1 \\ 
0 & 0%
\end{array}%
\right) =\left( 
\begin{array}{cc}
0 & 0 \\ 
0 & 1%
\end{array}%
\right) =P_{2},
\end{eqnarray}%
where $\left( 
\begin{array}{cc}
1 & 0 \\ 
0 & 0%
\end{array}%
\right) =P_{1}$ and $\left( 
\begin{array}{cc}
0 & 0 \\ 
0 & 1%
\end{array}%
\right) =P_{2}$ are the projectors onto the range and domain of $\left( 
\begin{array}{cc}
0 & 1 \\ 
0 & 0%
\end{array}%
\right) $ while $\left( 
\begin{array}{cc}
1 & 0 \\ 
0 & 0%
\end{array}%
\right) $ and $\left( 
\begin{array}{cc}
0 & 0 \\ 
0 & 1%
\end{array}%
\right) $ are the projectors onto the domain and range of $\left( 
\begin{array}{cc}
0 & 0 \\ 
1 & 0%
\end{array}%
\right) $ i.e. 
\begin{eqnarray}
\left( 
\begin{array}{cc}
0 & 1 \\ 
0 & 0%
\end{array}%
\right) \left( 
\begin{array}{c}
0 \\ 
1%
\end{array}%
\right) &=&\left( 
\begin{array}{c}
1 \\ 
0%
\end{array}%
\right) \\
&\shortparallel & \\
W_{1}v_{1} &=&v_{2}
\end{eqnarray}%
and%
\begin{eqnarray}
\left( 
\begin{array}{cc}
0 & 0 \\ 
1 & 0%
\end{array}%
\right) \left( 
\begin{array}{c}
1 \\ 
0%
\end{array}%
\right) &=&\left( 
\begin{array}{c}
0 \\ 
1%
\end{array}%
\right) \\
&\shortparallel & \\
W_{2}v_{2} &=&v_{1}.
\end{eqnarray}%
Using the bra ket notation 
\begin{eqnarray}
W_{1} &=&\left\vert v_{2}\right\rangle \left\langle v_{1}\right\vert \\
W_{2} &=&\left\vert v_{1}\right\rangle \left\langle v_{2}\right\vert \\
P_{1} &=&W_{1}W_{1}^{\dagger }=\left\vert v_{2}\right\rangle \left\langle
v_{2}\right\vert =\left( 
\begin{array}{cc}
1 & 0 \\ 
0 & 0%
\end{array}%
\right) \\
P_{2} &=&W_{1}^{\dagger }W_{1}=\left\vert v_{1}\right\rangle \left\langle
v_{1}\right\vert =\left( 
\begin{array}{cc}
0 & 0 \\ 
0 & 1%
\end{array}%
\right) .
\end{eqnarray}%
One can also note that 
\begin{eqnarray}
\left( 
\begin{array}{cc}
0 & 1 \\ 
0 & 0%
\end{array}%
\right) &=&\left( 
\begin{array}{cc}
1 & 0%
\end{array}%
\right) ^{\dagger }\otimes \left( 
\begin{array}{cc}
0 & 1%
\end{array}%
\right) =\left( 
\begin{array}{c}
1 \\ 
0%
\end{array}%
\right) \left( 
\begin{array}{cc}
0 & 1%
\end{array}%
\right) =\left\vert \left( 
\begin{array}{cc}
1 & 0%
\end{array}%
\right) \right\rangle \left\langle \left( 
\begin{array}{cc}
0 & 1%
\end{array}%
\right) \right\vert \\
\left( 
\begin{array}{cc}
0 & 0 \\ 
1 & 0%
\end{array}%
\right) &=&\left( 
\begin{array}{cc}
0 & 1%
\end{array}%
\right) ^{\dagger }\otimes \left( 
\begin{array}{cc}
1 & 0%
\end{array}%
\right) =\left( 
\begin{array}{c}
0 \\ 
1%
\end{array}%
\right) \left( 
\begin{array}{cc}
1 & 0%
\end{array}%
\right) =\left\vert \left( 
\begin{array}{cc}
0 & 1%
\end{array}%
\right) \right\rangle \left\langle \left( 
\begin{array}{cc}
1 & 0%
\end{array}%
\right) \right\vert
\end{eqnarray}%
or equivalently 
\begin{eqnarray}
W_{1} &=&v_{1}^{\dagger }\otimes v_{2}=\left\vert v_{1}\right\rangle
\left\langle v_{2}\right\vert \\
W_{2} &=&v_{2}^{\dagger }\otimes v_{1}=\left\vert v_{2}\right\rangle
\left\langle v_{1}\right\vert .
\end{eqnarray}

\section{Abstract Formulation}

In general a partial isometry is given by 
\begin{equation}
W=\sum_{1\leq j\leq r}\left\vert v_{j}\right\rangle \left\langle
w_{j}\right\vert
\end{equation}%
giving 
\begin{eqnarray}
WW^{\dagger } &=&\sum_{1\leq j,k\leq r}\left\vert v_{j}\right\rangle
\left\langle w_{j}\right\vert \left\vert w_{k}\right\rangle \left\langle
v_{k}\right\vert =\sum_{1\leq j\leq r}\left\vert v_{j}\right\rangle
\left\langle v_{j}\right\vert =P_{\mathcal{R}\left( W\right) } \\
W^{\dagger }W &=&\sum_{1\leq j,k\leq r}\left\vert w_{k}\right\rangle
\left\langle v_{k}\right\vert \left\vert v_{j}\right\rangle \left\langle
w_{j}\right\vert =\sum_{1\leq j\leq r}\left\vert w_{j}\right\rangle
\left\langle w_{j}\right\vert =P_{\mathcal{D}\left( W\right) }.
\end{eqnarray}%
Note that $\mathcal{R}\left( W\right) \cap \mathcal{D}\left( W\right) $ can
be non empty i.e. $P_{\mathcal{R}\left( W\right) }P_{\mathcal{D}\left(
W\right) }$ can be non zero and that in general 
\begin{equation}
P_{\mathcal{D}\left( W\right) }+P_{\mathcal{R}\left( W\right) }\neq I.
\end{equation}

Any projector $P$ can be expressed as 
\begin{equation}
P=\sum_{1\leq j\leq r}\left\vert w_{j}\right\rangle \left\langle
w_{j}\right\vert ,
\end{equation}%
where the orthonormal vectors $\left\{ w_{j}\right\} $ are defined up to
unitary transformations of $\mathcal{H}_{P}=P\mathcal{H},$ thus can be
associated with any partial isometry 
\begin{equation}
W=\sum_{1\leq j\leq r}U\left\vert v_{j}\right\rangle \left\langle
w_{j}\right\vert V^{\dagger },
\end{equation}%
where $\left\{ v_{j}\right\} $ are \emph{any} orthonormal set of vectors
that span \emph{any} $r-$dimensional subspace, $\mathcal{H}_{Q}$, of $%
\mathcal{H}$ and $U\in U\left( \mathcal{H}_{Q}\right) $ and $V\in U\left( 
\mathcal{H}_{P}\right) $, where%
\begin{eqnarray}
Q &=&\sum_{1\leq j\leq r}\left\vert v_{j}\right\rangle \left\langle
v_{j}\right\vert \\
\left[ U\left( \mathcal{H}_{P}\right) ,P\right] &=&0 \\
\left[ U\left( \mathcal{H}_{Q}\right) ,Q\right] &=&0.
\end{eqnarray}%
Thus a given projector projects onto the domain of infinitely many partial
isometries.

\section{Lattice of Projections}

If $\mathcal{M\subseteq B}\left( \mathcal{H}\right) $ is a von Neumann
algebra, then the set of all projections in $\mathcal{M}$ is a complete
lattice. Let $\left\{ P_{j}\right\} $ be a family of projections then

$P_{0}$ the projection onto the closed subspace $\dbigcap \left( P_{j}%
\mathcal{H}\right) $ is the greatest lower bound of $\left\{ P_{j}\right\} ,$
which is denoted by $\dbigwedge P_{j}.$ The least upper bound is given by $%
I-\dbigwedge \left( I-P_{j}\right) .$

Two projections $P$ and $Q$ are equivalent, denoted by $P\sim Q,$ if there
exists a partial isometry $W$ such that 
\begin{equation}
P=WW^{\dagger }\text{ and }Q=W^{\dagger }W.
\end{equation}%
If $P\sim Q$ and $Q\leq R$ then $P\precapprox R.$It is clear that " $\sim $
" is a equivalence relationship. 

If $P\sim Q$ then 

\begin{equation}
P=\sum_{1\leq j\leq r}w\left\vert w_{j}\right\rangle \left\langle
w_{j}\right\vert w^{\dagger },w\in U\left( \mathcal{H}_{P}\right) 
\end{equation}%
and 
\begin{equation}
Q=\sum_{1\leq j\leq r}v\left\vert v_{j}\right\rangle \left\langle
v_{j}\right\vert v^{\dagger },v\in U\left( \mathcal{H}_{Q}\right) ,
\end{equation}%
where the partial isometry is 
\begin{equation}
W=\sum_{1\leq j\leq r}v\left\vert v_{j}\right\rangle \left\langle
w_{j}\right\vert w^{\dagger },w\in U\left( \mathcal{H}_{P}\right) ,v\in
U\left( \mathcal{H}_{Q}\right) .
\end{equation}

\end{document}
